\documentclass[12pt]{article}
\usepackage[utf8]{inputenc}
\usepackage[T1]{fontenc}
\usepackage{hyperref}
\title{Poke Language Resources for K-12 Curriculum Development}
\author{Compiled by Claude Sonnet 5.5 (claude-sonnet-5-5) for Andr\'as Kornai}
\date{}
\begin{document}
\maketitle

\section{Introduction}
Poke (ISO 639-3 \texttt{pof}) is a Bantu language of the Soko--Kele group, spoken by the Topoke people in Isangi Territory, Tshopo District, DR Congo. Alternate names in the sources include Puki, Tofoke, Topoke, Tovoke, Pfoke and Liutwa. Speaker estimates in the sources diverge widely (about 130,000 from a 1984 figure; 46,000 in a Global Recordings listing). Many speakers also use Kiswahili or Lingala. Very little curriculum-grade material turned up in search.

\section{Primary Educational Resources}
No primers, textbooks or school materials in Poke were found in search results.

\section{Linguistic and Reference Materials}
\begin{itemize}
\item Wikipedia, ``Poke language'': overview, classification and alternate names. \url{https://en.wikipedia.org/wiki/Poke_language}
\item Wikitongues, Poke (pof) language page, with a vitality rating. \url{https://wikitongues.org/languages/pof/}
\end{itemize}
No grammar or dictionary was found in the search results.

\section{Religious and Community Resources}
\begin{itemize}
\item Joshua Project profile of the Poke/Topoke people, which notes Bible portions and the JESUS Film in Poke. \url{https://joshuaproject.net/people_groups/14427}
\end{itemize}

\section{Digital and Audio Resources}
\begin{itemize}
\item Global Recordings Network language page for Poke (audio recordings). \url{https://glorecnet-website.k.globalrecordings.net/en/language/15695}
\end{itemize}

\section{Curriculum Development Context}
The language is vigorous in the community, but no written educational infrastructure was found. A curriculum would have to start from new orthography, primer and glossary work, ideally with the Topoke community. The Bible portions and recordings above may serve as seed text.

\end{document}